本附录对应最终提交支撑包，内容分为支撑材料目录、完整文件清单、源程序说明和复现顺序。赛题附件 1 至附件 7 属于赛事统一提供的数据，按规范不重复放入支撑包；复现时应将原始目录“A题 (1)”放在支撑包根目录。附录 B 逐行收录建模使用的全部 11 个 Python 文件，不包含废弃模型、文档转换或临时绘图脚本。

\subsection*{A.1 支撑材料目录结构}

\begingroup
\small
\setlength{\tabcolsep}{4pt}
\begin{longtable}{p{0.31\textwidth}p{0.60\textwidth}}
\toprule
目录或文件 & 内容与用途 \\
\midrule
\endfirsthead
\toprule
目录或文件 & 内容与用途 \\
\midrule
\endhead
{\ttfamily\raggedright code/paper\_\allowbreak appendix/\par} & 四个问题、模型演进、最终后校准和官方指标复验所需的 11 个完整 Python 文件。 \\
\texttt{code/outputs/} & 问题三复跑时可选读取的整组留出失效诊断结果。 \\
\texttt{02\_关键结果数据/} & 基础积分、D0 校准、扰流特征、模型比较、最终提交和失效分析的必要结果。 \\
\texttt{03\_论文采用图表/} & 正文实际使用的 10 幅 PDF 图以及问题四工程图的可编辑 DrawIO 源文件。 \\
\texttt{04\_自主查阅资料/} & 文献理论索引与直接支撑积分、剖面修正、ANN、KRR 和物理约束的关键原文。 \\
\texttt{05\_复验与说明/} & 源码独立复跑记录、官方评价结果、文件哈希、AI 工具使用详情和终审报告。 \\
\texttt{requirements.txt} & Python 第三方依赖清单。 \\
\texttt{README.md} & 支撑包结构、原始附件放置方式和运行入口。 \\
\bottomrule
\end{longtable}
\endgroup

\subsection*{A.2 支撑材料完整文件清单}

下表逐项列出最终压缩包中的全部文件。文件哈希另见支撑包中的 \texttt{05\_复验与说明/文件清单与SHA256.csv}。

% Generated from the final support package. Do not edit by hand.
\begingroup
\scriptsize
\setlength{\tabcolsep}{3pt}
\renewcommand{\arraystretch}{0.92}
\begin{longtable}{p{0.055\textwidth}p{0.19\textwidth}p{0.665\textwidth}}
\toprule
序号 & 类别 & 相对路径 \\
\midrule
\endfirsthead
\toprule
序号 & 类别 & 相对路径 \\
\midrule
\endhead
1 & 关键结果数据 & {\ttfamily\raggedright 02\_\allowbreak{}关键结果数据/\allowbreak{}灵敏度与失效/\allowbreak{}krr\_\allowbreak{}leave\_\allowbreak{}one\_\allowbreak{}group\_\allowbreak{}failure\_\allowbreak{}modes.csv\par} \\
2 & 关键结果数据 & {\ttfamily\raggedright 02\_\allowbreak{}关键结果数据/\allowbreak{}灵敏度与失效/\allowbreak{}krr\_\allowbreak{}leave\_\allowbreak{}one\_\allowbreak{}group\_\allowbreak{}ood\_\allowbreak{}diagnostics.csv\par} \\
3 & 关键结果数据 & {\ttfamily\raggedright 02\_\allowbreak{}关键结果数据/\allowbreak{}灵敏度与失效/\allowbreak{}krr\_\allowbreak{}ood\_\allowbreak{}correlation\_\allowbreak{}sensitivity.csv\par} \\
4 & 关键结果数据 & {\ttfamily\raggedright 02\_\allowbreak{}关键结果数据/\allowbreak{}灵敏度与失效/\allowbreak{}krr\_\allowbreak{}pressure\_\allowbreak{}failure\_\allowbreak{}modes.csv\par} \\
5 & 关键结果数据 & {\ttfamily\raggedright 02\_\allowbreak{}关键结果数据/\allowbreak{}灵敏度与失效/\allowbreak{}krr\_\allowbreak{}pressure\_\allowbreak{}ood\_\allowbreak{}groups.csv\par} \\
6 & 关键结果数据 & {\ttfamily\raggedright 02\_\allowbreak{}关键结果数据/\allowbreak{}问题一\_\allowbreak{}基础积分/\allowbreak{}integration\_\allowbreak{}estimates\_\allowbreak{}long.csv\par} \\
7 & 关键结果数据 & {\ttfamily\raggedright 02\_\allowbreak{}关键结果数据/\allowbreak{}问题一\_\allowbreak{}基础积分/\allowbreak{}integration\_\allowbreak{}metrics\_\allowbreak{}by\_\allowbreak{}disturbed.csv\par} \\
8 & 关键结果数据 & {\ttfamily\raggedright 02\_\allowbreak{}关键结果数据/\allowbreak{}问题一\_\allowbreak{}基础积分/\allowbreak{}integration\_\allowbreak{}metrics\_\allowbreak{}by\_\allowbreak{}flow\_\allowbreak{}point.csv\par} \\
9 & 关键结果数据 & {\ttfamily\raggedright 02\_\allowbreak{}关键结果数据/\allowbreak{}问题一\_\allowbreak{}基础积分/\allowbreak{}integration\_\allowbreak{}overall\_\allowbreak{}metrics.csv\par} \\
10 & 关键结果数据 & {\ttfamily\raggedright 02\_\allowbreak{}关键结果数据/\allowbreak{}问题一\_\allowbreak{}基础积分/\allowbreak{}integration\_\allowbreak{}weights.csv\par} \\
11 & 关键结果数据 & {\ttfamily\raggedright 02\_\allowbreak{}关键结果数据/\allowbreak{}问题三\_\allowbreak{}特征与聚类/\allowbreak{}disturbance\_\allowbreak{}cluster\_\allowbreak{}cross\_\allowbreak{}disturbance.csv\par} \\
12 & 关键结果数据 & {\ttfamily\raggedright 02\_\allowbreak{}关键结果数据/\allowbreak{}问题三\_\allowbreak{}特征与聚类/\allowbreak{}disturbance\_\allowbreak{}cluster\_\allowbreak{}quality.csv\par} \\
13 & 关键结果数据 & {\ttfamily\raggedright 02\_\allowbreak{}关键结果数据/\allowbreak{}问题三\_\allowbreak{}特征与聚类/\allowbreak{}disturbance\_\allowbreak{}feature\_\allowbreak{}importance.csv\par} \\
14 & 关键结果数据 & {\ttfamily\raggedright 02\_\allowbreak{}关键结果数据/\allowbreak{}问题三\_\allowbreak{}特征与聚类/\allowbreak{}disturbance\_\allowbreak{}pca\_\allowbreak{}coordinates.csv\par} \\
15 & 关键结果数据 & {\ttfamily\raggedright 02\_\allowbreak{}关键结果数据/\allowbreak{}问题三\_\allowbreak{}特征与聚类/\allowbreak{}disturbance\_\allowbreak{}pca\_\allowbreak{}explained\_\allowbreak{}variance.csv\par} \\
16 & 关键结果数据 & {\ttfamily\raggedright 02\_\allowbreak{}关键结果数据/\allowbreak{}问题二\_\allowbreak{}D0校准/\allowbreak{}d0\_\allowbreak{}model\_\allowbreak{}by\_\allowbreak{}flow\_\allowbreak{}d0\_\allowbreak{}only.csv\par} \\
17 & 关键结果数据 & {\ttfamily\raggedright 02\_\allowbreak{}关键结果数据/\allowbreak{}问题二\_\allowbreak{}D0校准/\allowbreak{}d0\_\allowbreak{}model\_\allowbreak{}overall\_\allowbreak{}d0\_\allowbreak{}only.csv\par} \\
18 & 关键结果数据 & {\ttfamily\raggedright 02\_\allowbreak{}关键结果数据/\allowbreak{}问题二\_\allowbreak{}D0校准/\allowbreak{}phys6\_\allowbreak{}d0\_\allowbreak{}curve\_\allowbreak{}coefficients\_\allowbreak{}main\_\allowbreak{}split.csv\par} \\
19 & 关键结果数据 & {\ttfamily\raggedright 02\_\allowbreak{}关键结果数据/\allowbreak{}问题四\_\allowbreak{}最终模型/\allowbreak{}d0\_\allowbreak{}rbf\_\allowbreak{}postcalibrated\_\allowbreak{}dual\_\allowbreak{}kernel\_\allowbreak{}krr\_\allowbreak{}candidate\_\allowbreak{}ranking.csv\par} \\
20 & 关键结果数据 & {\ttfamily\raggedright 02\_\allowbreak{}关键结果数据/\allowbreak{}问题四\_\allowbreak{}最终模型/\allowbreak{}d0\_\allowbreak{}rbf\_\allowbreak{}postcalibrated\_\allowbreak{}dual\_\allowbreak{}kernel\_\allowbreak{}krr\_\allowbreak{}evaluation\_\allowbreak{}groups.csv\par} \\
21 & 关键结果数据 & {\ttfamily\raggedright 02\_\allowbreak{}关键结果数据/\allowbreak{}问题四\_\allowbreak{}最终模型/\allowbreak{}d0\_\allowbreak{}rbf\_\allowbreak{}postcalibrated\_\allowbreak{}dual\_\allowbreak{}kernel\_\allowbreak{}krr\_\allowbreak{}evaluation\_\allowbreak{}summary.csv\par} \\
22 & 关键结果数据 & {\ttfamily\raggedright 02\_\allowbreak{}关键结果数据/\allowbreak{}问题四\_\allowbreak{}最终模型/\allowbreak{}d0\_\allowbreak{}rbf\_\allowbreak{}postcalibrated\_\allowbreak{}dual\_\allowbreak{}kernel\_\allowbreak{}krr\_\allowbreak{}evaluation\_\allowbreak{}windows.csv\par} \\
23 & 关键结果数据 & {\ttfamily\raggedright 02\_\allowbreak{}关键结果数据/\allowbreak{}问题四\_\allowbreak{}最终模型/\allowbreak{}d0\_\allowbreak{}rbf\_\allowbreak{}postcalibrated\_\allowbreak{}dual\_\allowbreak{}kernel\_\allowbreak{}krr\_\allowbreak{}model.npz\par} \\
24 & 关键结果数据 & {\ttfamily\raggedright 02\_\allowbreak{}关键结果数据/\allowbreak{}问题四\_\allowbreak{}最终模型/\allowbreak{}d0\_\allowbreak{}rbf\_\allowbreak{}postcalibrated\_\allowbreak{}dual\_\allowbreak{}kernel\_\allowbreak{}krr\_\allowbreak{}model\_\allowbreak{}meta.json\par} \\
25 & 关键结果数据 & {\ttfamily\raggedright 02\_\allowbreak{}关键结果数据/\allowbreak{}问题四\_\allowbreak{}最终模型/\allowbreak{}submission\_\allowbreak{}final.csv\par} \\
26 & 关键结果数据 & {\ttfamily\raggedright 02\_\allowbreak{}关键结果数据/\allowbreak{}问题四\_\allowbreak{}模型比较/\allowbreak{}ann\_\allowbreak{}krr\_\allowbreak{}same\_\allowbreak{}split\_\allowbreak{}assignments.csv\par} \\
27 & 关键结果数据 & {\ttfamily\raggedright 02\_\allowbreak{}关键结果数据/\allowbreak{}问题四\_\allowbreak{}模型比较/\allowbreak{}ann\_\allowbreak{}krr\_\allowbreak{}same\_\allowbreak{}split\_\allowbreak{}comparison.csv\par} \\
28 & 关键结果数据 & {\ttfamily\raggedright 02\_\allowbreak{}关键结果数据/\allowbreak{}问题四\_\allowbreak{}模型比较/\allowbreak{}final\_\allowbreak{}model\_\allowbreak{}evolution\_\allowbreak{}comparison.csv\par} \\
29 & 关键结果数据 & {\ttfamily\raggedright 02\_\allowbreak{}关键结果数据/\allowbreak{}问题四\_\allowbreak{}模型比较/\allowbreak{}fullfit\_\allowbreak{}interpolation\_\allowbreak{}evaluation\_\allowbreak{}summary.csv\par} \\
30 & 关键结果数据 & {\ttfamily\raggedright 02\_\allowbreak{}关键结果数据/\allowbreak{}问题四\_\allowbreak{}模型比较/\allowbreak{}krr\_\allowbreak{}champion\_\allowbreak{}multisplit\_\allowbreak{}metrics.csv\par} \\
31 & 关键结果数据 & {\ttfamily\raggedright 02\_\allowbreak{}关键结果数据/\allowbreak{}问题四\_\allowbreak{}模型比较/\allowbreak{}krr\_\allowbreak{}champion\_\allowbreak{}multisplit\_\allowbreak{}pass\_\allowbreak{}rates.csv\par} \\
32 & 关键结果数据 & {\ttfamily\raggedright 02\_\allowbreak{}关键结果数据/\allowbreak{}问题四\_\allowbreak{}模型比较/\allowbreak{}krr\_\allowbreak{}champion\_\allowbreak{}multisplit\_\allowbreak{}summary.csv\par} \\
33 & 关键结果数据 & {\ttfamily\raggedright 02\_\allowbreak{}关键结果数据/\allowbreak{}问题四\_\allowbreak{}模型比较/\allowbreak{}krr\_\allowbreak{}mainline\_\allowbreak{}ablation\_\allowbreak{}main\_\allowbreak{}validation.csv\par} \\
34 & 关键结果数据 & {\ttfamily\raggedright 02\_\allowbreak{}关键结果数据/\allowbreak{}问题四\_\allowbreak{}模型比较/\allowbreak{}stratified\_\allowbreak{}model\_\allowbreak{}search\_\allowbreak{}best\_\allowbreak{}split.csv\par} \\
35 & 关键结果数据 & {\ttfamily\raggedright 02\_\allowbreak{}关键结果数据/\allowbreak{}问题四\_\allowbreak{}模型比较/\allowbreak{}stratified\_\allowbreak{}model\_\allowbreak{}search\_\allowbreak{}multisplit\_\allowbreak{}summary.csv\par} \\
36 & 关键结果数据 & {\ttfamily\raggedright 02\_\allowbreak{}关键结果数据/\allowbreak{}问题四\_\allowbreak{}模型比较/\allowbreak{}stratified\_\allowbreak{}model\_\allowbreak{}search\_\allowbreak{}pressure\_\allowbreak{}split.csv\par} \\
37 & 论文采用图表 & {\ttfamily\raggedright 03\_\allowbreak{}论文采用图表/\allowbreak{}fig\_\allowbreak{}disturbance\_\allowbreak{}feature\_\allowbreak{}importance.pdf\par} \\
38 & 论文采用图表 & {\ttfamily\raggedright 03\_\allowbreak{}论文采用图表/\allowbreak{}fig\_\allowbreak{}disturbance\_\allowbreak{}pca\_\allowbreak{}clusters.pdf\par} \\
39 & 论文采用图表 & {\ttfamily\raggedright 03\_\allowbreak{}论文采用图表/\allowbreak{}fig\_\allowbreak{}disturbance\_\allowbreak{}profile\_\allowbreak{}representativeness.pdf\par} \\
40 & 论文采用图表 & {\ttfamily\raggedright 03\_\allowbreak{}论文采用图表/\allowbreak{}fig\_\allowbreak{}integration\_\allowbreak{}flow\_\allowbreak{}mae.pdf\par} \\
41 & 论文采用图表 & {\ttfamily\raggedright 03\_\allowbreak{}论文采用图表/\allowbreak{}fig\_\allowbreak{}krr\_\allowbreak{}ablation\_\allowbreak{}metrics.pdf\par} \\
42 & 论文采用图表 & {\ttfamily\raggedright 03\_\allowbreak{}论文采用图表/\allowbreak{}fig\_\allowbreak{}krr\_\allowbreak{}ood\_\allowbreak{}vs\_\allowbreak{}error.pdf\par} \\
43 & 论文采用图表 & {\ttfamily\raggedright 03\_\allowbreak{}论文采用图表/\allowbreak{}fig\_\allowbreak{}leave\_\allowbreak{}one\_\allowbreak{}group\_\allowbreak{}error\_\allowbreak{}heatmap.pdf\par} \\
44 & 论文采用图表 & {\ttfamily\raggedright 03\_\allowbreak{}论文采用图表/\allowbreak{}fig\_\allowbreak{}model\_\allowbreak{}mae\_\allowbreak{}comparison.pdf\par} \\
45 & 论文采用图表 & {\ttfamily\raggedright 03\_\allowbreak{}论文采用图表/\allowbreak{}fig\_\allowbreak{}prediction\_\allowbreak{}consistency.pdf\par} \\
46 & 论文采用图表 & {\ttfamily\raggedright 03\_\allowbreak{}论文采用图表/\allowbreak{}fig\_\allowbreak{}problem4\_\allowbreak{}engineering\_\allowbreak{}pipeline.drawio\par} \\
47 & 论文采用图表 & {\ttfamily\raggedright 03\_\allowbreak{}论文采用图表/\allowbreak{}fig\_\allowbreak{}problem4\_\allowbreak{}engineering\_\allowbreak{}pipeline.pdf\par} \\
48 & 自主查阅资料 & {\ttfamily\raggedright 04\_\allowbreak{}自主查阅资料/\allowbreak{}关键原文/\allowbreak{}ANN\_\allowbreak{}data\_\allowbreak{}integration\_\allowbreak{}for\_\allowbreak{}multipath\_\allowbreak{}ultrasonic\_\allowbreak{}flowmeter.pdf\par} \\
49 & 自主查阅资料 & {\ttfamily\raggedright 04\_\allowbreak{}自主查阅资料/\allowbreak{}关键原文/\allowbreak{}Calculation\_\allowbreak{}geometric\_\allowbreak{}flow\_\allowbreak{}profile\_\allowbreak{}correction\_\allowbreak{}factor.pdf\par} \\
50 & 自主查阅资料 & {\ttfamily\raggedright 04\_\allowbreak{}自主查阅资料/\allowbreak{}关键原文/\allowbreak{}Flow\_\allowbreak{}profile\_\allowbreak{}identification\_\allowbreak{}by\_\allowbreak{}multi\_\allowbreak{}channel\_\allowbreak{}measurements.pdf\par} \\
51 & 自主查阅资料 & {\ttfamily\raggedright 04\_\allowbreak{}自主查阅资料/\allowbreak{}关键原文/\allowbreak{}Generalized\_\allowbreak{}velocity\_\allowbreak{}profile\_\allowbreak{}evaluation.pdf\par} \\
52 & 自主查阅资料 & {\ttfamily\raggedright 04\_\allowbreak{}自主查阅资料/\allowbreak{}关键原文/\allowbreak{}In\_\allowbreak{}use\_\allowbreak{}measurement\_\allowbreak{}based\_\allowbreak{}on\_\allowbreak{}machine\_\allowbreak{}learning.pdf\par} \\
53 & 自主查阅资料 & {\ttfamily\raggedright 04\_\allowbreak{}自主查阅资料/\allowbreak{}关键原文/\allowbreak{}Physically\_\allowbreak{}constrained\_\allowbreak{}GAN\_\allowbreak{}data\_\allowbreak{}augmentation.pdf\par} \\
54 & 自主查阅资料 & {\ttfamily\raggedright 04\_\allowbreak{}自主查阅资料/\allowbreak{}理论索引/\allowbreak{}LITERATURE\_\allowbreak{}BASIS.md\par} \\
55 & 自主查阅资料 & {\ttfamily\raggedright 04\_\allowbreak{}自主查阅资料/\allowbreak{}理论索引/\allowbreak{}LITERATURE\_\allowbreak{}KEYPOINTS\_\allowbreak{}BOOKLET.md\par} \\
56 & 自主查阅资料 & {\ttfamily\raggedright 04\_\allowbreak{}自主查阅资料/\allowbreak{}理论索引/\allowbreak{}LITERATURE\_\allowbreak{}KEYPOINTS\_\allowbreak{}BOOKLET.pdf\par} \\
57 & 自主查阅资料 & {\ttfamily\raggedright 04\_\allowbreak{}自主查阅资料/\allowbreak{}理论索引/\allowbreak{}LITERATURE\_\allowbreak{}LIBRARY.md\par} \\
58 & 自主查阅资料 & {\ttfamily\raggedright 04\_\allowbreak{}自主查阅资料/\allowbreak{}理论索引/\allowbreak{}references.tex\par} \\
59 & 自主查阅资料 & {\ttfamily\raggedright 04\_\allowbreak{}自主查阅资料/\allowbreak{}资料来源与用途说明.md\par} \\
60 & 复验与说明 & {\ttfamily\raggedright 05\_\allowbreak{}复验与说明/\allowbreak{}APPENDIX\_\allowbreak{}CODE\_\allowbreak{}TRUTH\_\allowbreak{}AUDIT.md\par} \\
61 & 复验与说明 & {\ttfamily\raggedright 05\_\allowbreak{}复验与说明/\allowbreak{}package\_\allowbreak{}validation.json\par} \\
62 & 复验与说明 & {\ttfamily\raggedright 05\_\allowbreak{}复验与说明/\allowbreak{}PYCHARM\_\allowbreak{}SETUP.md\par} \\
63 & 复验与说明 & {\ttfamily\raggedright 05\_\allowbreak{}复验与说明/\allowbreak{}图表数据与支撑包终审报告.md\par} \\
64 & 复验与说明 & {\ttfamily\raggedright 05\_\allowbreak{}复验与说明/\allowbreak{}官方脚本最终复验/\allowbreak{}evaluation\_\allowbreak{}groups.csv\par} \\
65 & 复验与说明 & {\ttfamily\raggedright 05\_\allowbreak{}复验与说明/\allowbreak{}官方脚本最终复验/\allowbreak{}evaluation\_\allowbreak{}summary.csv\par} \\
66 & 复验与说明 & {\ttfamily\raggedright 05\_\allowbreak{}复验与说明/\allowbreak{}官方脚本最终复验/\allowbreak{}evaluation\_\allowbreak{}windows.csv\par} \\
67 & 复验与说明 & {\ttfamily\raggedright 05\_\allowbreak{}复验与说明/\allowbreak{}文件清单与SHA256.csv\par} \\
68 & 复验与说明 & {\ttfamily\raggedright 05\_\allowbreak{}复验与说明/\allowbreak{}独立源码复跑状态.csv\par} \\
69 & 根目录说明 & {\ttfamily\raggedright AI工具使用详情.pdf\par} \\
70 & 源程序与运行依赖 & {\ttfamily\raggedright code/\allowbreak{}outputs/\allowbreak{}krr\_\allowbreak{}leave\_\allowbreak{}one\_\allowbreak{}group\_\allowbreak{}failure\_\allowbreak{}modes.csv\par} \\
71 & 源程序与运行依赖 & {\ttfamily\raggedright code/\allowbreak{}paper\_\allowbreak{}appendix/\allowbreak{}ann\_\allowbreak{}krr\_\allowbreak{}same\_\allowbreak{}split.py\par} \\
72 & 源程序与运行依赖 & {\ttfamily\raggedright code/\allowbreak{}paper\_\allowbreak{}appendix/\allowbreak{}compact\_\allowbreak{}rbf\_\allowbreak{}krr.py\par} \\
73 & 源程序与运行依赖 & {\ttfamily\raggedright code/\allowbreak{}paper\_\allowbreak{}appendix/\allowbreak{}final\_\allowbreak{}postcalibrated\_\allowbreak{}dual\_\allowbreak{}kernel\_\allowbreak{}krr.py\par} \\
74 & 源程序与运行依赖 & {\ttfamily\raggedright code/\allowbreak{}paper\_\allowbreak{}appendix/\allowbreak{}fullfit\_\allowbreak{}expression\_\allowbreak{}baseline.py\par} \\
75 & 源程序与运行依赖 & {\ttfamily\raggedright code/\allowbreak{}paper\_\allowbreak{}appendix/\allowbreak{}model\_\allowbreak{}evolution\_\allowbreak{}ridge\_\allowbreak{}ann.py\par} \\
76 & 源程序与运行依赖 & {\ttfamily\raggedright code/\allowbreak{}paper\_\allowbreak{}appendix/\allowbreak{}official\_\allowbreak{}evaluator.py\par} \\
77 & 源程序与运行依赖 & {\ttfamily\raggedright code/\allowbreak{}paper\_\allowbreak{}appendix/\allowbreak{}physical\_\allowbreak{}feature\_\allowbreak{}library.py\par} \\
78 & 源程序与运行依赖 & {\ttfamily\raggedright code/\allowbreak{}paper\_\allowbreak{}appendix/\allowbreak{}physical\_\allowbreak{}stability\_\allowbreak{}dual\_\allowbreak{}kernel\_\allowbreak{}krr.py\par} \\
79 & 源程序与运行依赖 & {\ttfamily\raggedright code/\allowbreak{}paper\_\allowbreak{}appendix/\allowbreak{}problem1\_\allowbreak{}multichannel\_\allowbreak{}integration.py\par} \\
80 & 源程序与运行依赖 & {\ttfamily\raggedright code/\allowbreak{}paper\_\allowbreak{}appendix/\allowbreak{}problem2\_\allowbreak{}d0\_\allowbreak{}calibration.py\par} \\
81 & 源程序与运行依赖 & {\ttfamily\raggedright code/\allowbreak{}paper\_\allowbreak{}appendix/\allowbreak{}problem3\_\allowbreak{}disturbance\_\allowbreak{}features.py\par} \\
82 & 源程序与运行依赖 & {\ttfamily\raggedright code/\allowbreak{}paper\_\allowbreak{}appendix/\allowbreak{}README.md\par} \\
83 & 根目录说明 & {\ttfamily\raggedright README.md\par} \\
84 & 根目录说明 & {\ttfamily\raggedright requirements.txt\par} \\
\bottomrule
\end{longtable}
\endgroup

% support_file_count=84


\subsection*{A.3 完整源程序及其作用}

\begingroup
\scriptsize
\setlength{\tabcolsep}{3pt}
\renewcommand{\arraystretch}{1.08}
\begin{longtable}{p{0.39\textwidth}p{0.52\textwidth}}
\toprule
源程序 & 作用 \\
\midrule
\endfirsthead
\toprule
源程序 & 作用 \\
\midrule
\endhead
{\ttfamily\raggedright problem1\_multichannel\_\allowbreak integration.py\par} & 重建五声道观测，计算 Phys6、OWICS、Lagrange、等权和单声道估计。 \\
{\ttfamily\raggedright problem2\_d0\_\allowbreak calibration.py\par} & 在 D0 无扰流窗口上建立正常工况校准与多声道权重模型。 \\
{\ttfamily\raggedright model\_evolution\_\allowbreak ridge\_ann.py\par} & 公共数据读取、评价指标、Ridge 与 ANN 残差模型及模型演进对照。 \\
{\ttfamily\raggedright physical\_feature\_\allowbreak library.py\par} & 构造紧凑剖面特征和运行时物理特征。 \\
{\ttfamily\raggedright compact\_rbf\_krr.py\par} & 实现标签盲分层、RBF-KRR、超参数搜索和紧凑模型评价。 \\
{\ttfamily\raggedright ann\_krr\_same\_split.py\par} & 在同一 $118/41$ 划分上比较 ANN 与紧凑 RBF-KRR。 \\
{\ttfamily\raggedright problem3\_disturbance\_\allowbreak features.py\par} & 计算扰流特征重要性、PCA、聚类及其与失效类型的关系。 \\
{\ttfamily\raggedright fullfit\_expression\_\allowbreak baseline.py\par} & 建立高维全量近插值对照，用于样本内表达上限检验。 \\
{\ttfamily\raggedright physical\_stability\_\allowbreak dual\_kernel\_krr.py\par} & 从附件 6 提取逐秒稳定性特征，并融合物理核与稳定性核。 \\
{\ttfamily\raggedright final\_postcalibrated\_\allowbreak dual\_kernel\_krr.py\par} & 在冻结双核 KRR 上搜索受约束 D0-RBF 后校准并生成最终提交。 \\
{\ttfamily\raggedright official\_evaluator.py\par} & 按题设口径计算相对误差、评价组、$u_{\mathrm{nor},L}$、$u_{\mathrm{nor},r}$ 和 $u_{\mathrm{nor},d}$。 \\
\bottomrule
\end{longtable}
\endgroup

上述文件构成一条闭合依赖链。公共函数只在这 11 个文件之间导入，运行时不需要旧版候选模型代码。各脚本会自动创建独立的结果、图形和报告目录，因此可在新目录中直接复跑，不会覆盖论文采用的验收结果。

\subsection*{A.4 运行顺序与官方复验}

\begin{lstlisting}[language=bash,basicstyle=\ttfamily\scriptsize,aboveskip=0.25em,belowskip=0pt]
python code/paper_appendix/problem1_multichannel_integration.py
python code/paper_appendix/problem2_d0_calibration.py
python code/paper_appendix/model_evolution_ridge_ann.py
python code/paper_appendix/physical_feature_library.py
python code/paper_appendix/compact_rbf_krr.py
python code/paper_appendix/ann_krr_same_split.py
python code/paper_appendix/problem3_disturbance_features.py
python code/paper_appendix/fullfit_expression_baseline.py
python code/paper_appendix/physical_stability_dual_kernel_krr.py
python code/paper_appendix/final_postcalibrated_dual_kernel_krr.py
python code/paper_appendix/official_evaluator.py \
  code/paper_appendix/verification_outputs/submission_final.csv \
  --output-dir code/paper_appendix/official_evaluation
\end{lstlisting}

独立干净目录复跑时，上述 11 个程序退出码均为 0。最后一条命令重新评价生成的 \texttt{submission\_final.csv}，得到 159 个窗口、30/30 评价组通过，$u_{\mathrm{nor},L}=0.009585\%$、$u_{\mathrm{nor},r}=0.033630\%$、$u_{\mathrm{nor},d}=0.048435\%$。该结果是同样本全量标定证据；标签盲划分仍单独用于讨论泛化能力，两种口径不得混用。
